\question [20] Na pista de testes de uma montadora de automóveis, foram feitas medições do comprimento da pista e do tempo gasto por um certo veículo para percorrê-la. Os valores obtidos foram, respectivamente, \SI{1030}{\meter} e \SI{25}{\second}. Levando-se em conta a precisão das medidas efetuadas, é correto afirmar que a velocidade média desenvolvida pelo citado veículo foi, em \si[per-mode=fraction]{\meter\per\second}, de  


		\begin{EnvUplevel}
			\begin{choices}
			
			\choice 45,3  
			\choice 47     
			\correctchoice 41,2       
			\choice 53      
			\choice 40  
				
			\end{choices}
		\end{EnvUplevel}
		
		\begin{EnvUplevel}
		\begin{solution}
		Foi dado pelo enunciado que a distância percorrida \Delta S foi 1030 metros  metros e a duração do movimento  \Delta T foi de 25 segundos. Logo,
		$$
		V_m=\frac{\Delta S}{\Delta T}=\frac{1030}{25}
		$$$$
		V_m=\SI[per-mode=fraction]{41,2}{\meter\per\second}
		$$
		\end{solution}
		\end{EnvUplevel}


	
